\chapter{Fondamenti e algoritmi}\label{ch:teoria}
\section{Dal recommender al bandit contestuale}
Nel Multi-Armed Bandit classico una policy sceglie un'azione $a_t$ e riceve la sola ricompensa associata alla scelta. L'esplorazione acquisisce informazione su alternative incerte; l'exploitation sfrutta le stime correnti. UCB aggiunge alla media empirica un bonus di incertezza:
\begin{equation}
 U_a(t)=\widehat{\mu}_a(t)+\sqrt{\frac{2\ln t}{n_a(t)}},
\end{equation}
per azioni con $n_a(t)>0$, nel caso classico di reward indipendenti, stazionari e limitati in $[0,1]$ per ciascun braccio \cite{auer2002}. UCB1 inizializza provando ogni azione almeno una volta; equivalentemente, un'azione non osservata ha priorità esplorativa. La formula fornisce soltanto l'intuizione: il progetto implementa LinUCB, non UCB1.

Nel CMAB il vettore $x_{u,a,t}\in\mathbb{R}^{d}$ riassume utente, item e stato prima della decisione. Generalizzare fra candidati è utile quando non tutti gli item hanno feedback individuali sufficienti. Qui $d=15$ e $\mathcal A_u(t)$ varia; il ritardo separa decisioni e aggiornamenti disponibili.

\section{LinUCB effettivamente implementato}
Il modello lineare condiviso usa una matrice $A\in\mathbb{R}^{15\times15}$ e un vettore $b\in\mathbb{R}^{15}$, inizializzati con $A=I$ e $b=0$. Non sono mantenute matrici separate per ciascuna traccia o ciascun utente. Date le sole azioni con feedback maturo, valgono
\begin{align}
A_t &= I+\sum_{\tau:\,\tau+\Delta\leq t}x_\tau x_\tau^\top,
&b_t &= \sum_{\tau:\,\tau+\Delta\leq t}r_\tau x_\tau,\\
\widehat{\theta}_t&=A_t^{-1}b_t,
&s(x)&=x^\top\widehat{\theta}_t+
\alpha\sqrt{x^\top A_t^{-1}x}.\label{eq:linucb}
\end{align}
Il vettore $x_\tau$ è quello standardizzato al momento della decisione, conservato nel feedback. Il parametro $\alpha$ controlla il bonus; $\alpha=0$ conserva soltanto la stima lineare, senza diventare la baseline Personalized Greedy. La regola deriva dall'approccio LinUCB \cite{li2010}, adattato qui a un modello condiviso sulle feature delle azioni.

Il codice usa Cholesky e sostituzioni triangolari senza costruire l'inversa. La selezione non modifica $A,b$; solo il feedback maturo della Top-1 applica $A\leftarrow A+xx^\top$, $b\leftarrow b+rx$. Non esiste un'intercetta esplicita: manca un livello di base liberamente stimabile e, per $x=0$, la componente media è zero. La personalizzazione passa dal contesto, con coefficienti condivisi. Il reward logaritmico non è limitato a $[0,1]$: non si rivendicano garanzie di regret derivate per altre distribuzioni.

\section{Baseline e interpretazione del confronto}
Random, permutazione uniforme con generatore isolato, rappresenta l'esplorazione non personalizzata. Global Popularity usa gli ascolti globali dell'artista come riferimento aggregato. Personalized Greedy sfrutta invece le preferenze storiche mediante il conteggio utente--artista, con tie-break per popolarità della traccia e chiave canonica. Il confronto distingue così il valore della personalizzazione da quello dei segnali globali, entro lo stesso pool.

$\varepsilon$-Greedy introduce una dose controllata di esplorazione casuale: con probabilità $\varepsilon$ permuta il pool, altrimenti usa l'intero ranking Greedy. Il ramo è estratto una volta per decisione e riguarda anche la Top-K. Non apprende medie di reward. LinUCB guida invece le scelte con stima contestuale e incertezza, studiando un'alternativa al distacco casuale dalle preferenze storiche.

Le strategie coprono diversi compromessi exploration--exploitation, ma LinUCB cambia anche rappresentazione e apprendimento: senza ablation il confronto non isola causalmente il solo bonus UCB.
